\chapter*{Preface}
\addcontentsline{toc}{chapter}{Preface}

Cycle III closed by asking what compositional relations organize RSVP structures, and answered with functors, natural transformations, adjunctions, limits, colimits, and finally the monoidal closure Book 15 constructed. The present book asks a different question: when do locally specified RSVP structures constitute a coherent global object at all? This is not a further compositional question but a question of coherence emerging from locality, and it inaugurates the sheaf-theoretic phase of RSVP theory directly. An entropy sheaf $\mathcal{S}$ assigns to every open region $U \subset M$ the space of admissible entropy fields $\SEntropy|_U$ satisfying local smoothness and conservation, together with restriction maps encoding boundary compatibility, and this formalism captures the essential feature of RSVP cosmology that Book 9's Chapter 10 and Book 14's Chapter 5 both borrowed ahead of their own place in the series to state: coherence is never imposed globally but emerges from consistent local smoothing.

\chapter*{Diagrammatic Overview}
\addcontentsline{toc}{chapter}{Diagrammatic Overview}

The relationship between the monoidal composition of Book 15 and the sheaf gluing of this book is one of scale rather than opposition: Book 15's tensor product combined whole objects of $\mathbf{RSVPField}$ with one another, while the present book's sheaf gluing combines a single object's own local pieces into that object itself. Composition, in the sense this cycle has developed it since Book 11, happens both between systems and within them, and this book supplies the within.

\chapter*{Reading Note}
\addcontentsline{toc}{chapter}{Reading Note}

Every act of understanding is the gluing of nearby smoothness. This book asks the reader to take that claim as literally as the categorical machinery it develops permits: understanding, throughout these twelve chapters, just is the existence of a global section over a covering of local ones, no further gloss required.

\part{Sheaf Foundations for Thermodynamic Fields}

\chapter{The Topological Manifold of Entropy}

The base space $M$ of this book's construction carries a topology generated by regions of homogeneous curvature, with open sets $U_i$ representing neighborhoods of approximate equilibrium and morphisms between them given by ordinary inclusion maps preserving boundary entropy. This is the identical base manifold Book 9's Chapter 10 and Book 14's Chapter 5 both already assumed when they borrowed this book's vocabulary ahead of its own introduction, and the present chapter makes explicit what those earlier chapters could only state informally: $M$ is an ordinary topological space, its open sets are ordinary opens, and nothing about the construction that follows requires any generalization of that ordinary notion of covering.

The physical interpretation offered here treats these open sets as overlapping thermodynamic patches of the plenum, regions within which the entropy field's local behavior can be described independently of what occurs elsewhere, with the overlaps between such patches being exactly where the gluing machinery of Part II earns its keep.

